% Archive record 67, line 1.
\hypertarget{part-4-67-0001-iv-d-returns-sharpe-ratio-performance-evaluation}{%
\section{IV-D --- Returns, Sharpe Ratio \& Performance Evaluation}\label{part-4-67-0001-iv-d-returns-sharpe-ratio-performance-evaluation}}

上一阶段 IV-C，我们研究了 Financial Factors、Prediction Models、Information Coefficient 和 Rank IC。

但现在出现了一个非常重要的问题：

\textbf{假设 RSI Agent 不断提高金融预测模型的 IC，那么它所创造的经济价值是否也会不断提高？}

答案是不一定。

在金融研究中，我们需要严格区分三个不同层级的性能：

\begin{readertable}
\begin{longtable}[]{@{}
  >{\raggedright\arraybackslash}p{(\columnwidth - 4\tabcolsep) * \real{0.28}}
  >{\raggedright\arraybackslash}p{(\columnwidth - 4\tabcolsep) * \real{0.35}}
  >{\raggedright\arraybackslash}p{(\columnwidth - 4\tabcolsep) * \real{0.34}}@{}}
\toprule
Performance Level & Research Question & Typical Metrics \\
\midrule
\endhead
Statistical Performance & 预测是否准确？ & MSE、MAE、\(R^2\) \\
Ranking Performance & 是否正确识别相对表现？ & IC、Rank IC、ICIR \\
Economic Performance & 预测转化为投资组合后是否有效？ & Return、Sharpe Ratio、Drawdown \\
\bottomrule
\end{longtable}
\end{readertable}

这三类性能之间存在联系，但并不等价。

对于 RSI 而言，选择哪一种指标作为 Objective Function，将直接影响 Agent 的搜索方向、Revision 类型以及潜在的过拟合机制。

\hypertarget{part-4-67-0001-learning-roadmap}{%
\subsection{0. Learning Roadmap}\label{part-4-67-0001-learning-roadmap}}

\begin{readertable}
\begin{longtable}[]{@{}
  >{\raggedright\arraybackslash}p{(\columnwidth - 4\tabcolsep) * \real{0.19}}
  >{\raggedright\arraybackslash}p{(\columnwidth - 4\tabcolsep) * \real{0.55}}
  >{\raggedright\arraybackslash}p{(\columnwidth - 4\tabcolsep) * \real{0.21}}@{}}
\toprule
Section & Topic & Priority \\
\midrule
\endhead
IV-A & Financial Time Series \& DGP & \\
IV-B & Low SNR \& Predictability & \\
IV-C & Financial Factors \& ML Models & \\
\textbf{IV-D} & \textbf{Returns, Sharpe Ratio \& Evaluation} & 当前 \\
IV-E & Backtest Overfitting, PBO \& DSR & 待学习 \\
IV-F & Market Regime \& Robustness & 待学习 \\
\bottomrule
\end{longtable}
\end{readertable}

本阶段最重要的新数学内容是 \textbf{Sharpe Ratio 的统计性质}。我们将利用 Linear Algebra、Conditional Expectation、CLT 和 Delta Method 解释为什么这个广泛使用的金融指标可能被 RSI 系统严重高估。

\hypertarget{part-4-67-0001-from-prediction-to-portfolioux4eceux9884ux6d4bux5230ux6295ux8d44ux7ec4ux5408}{%
\subsection{1. From Prediction to Portfolio（从预测到投资组合）}\label{part-4-67-0001-from-prediction-to-portfolioux4eceux9884ux6d4bux5230ux6295ux8d44ux7ec4ux5408}}

\hypertarget{part-4-67-0001-prediction-is-not-a-trading-decision}{%
\subsubsection{1.1 Prediction Is Not a Trading Decision}\label{part-4-67-0001-prediction-is-not-a-trading-decision}}

沿用 IV-C 中的记号：
\[\hat r_{i,t+1}=h(X_{i,t}).\]

这是资产 \(i\) 在时刻 \(t\) 对下一期收益率的预测。

但是，预测模型本身并没有决定投资组合如何分配资金。我们还需要一个 Portfolio Construction Function。

定义：
\[\boldsymbol w_t
=
\psi(\hat{\boldsymbol r}_{t+1}),\]

其中：
\[\boldsymbol w_t=
\begin{pmatrix}
w_{1,t}\\
\vdots\\
w_{N,t}
\end{pmatrix}\]

表示投资组合中各资产的期初权重，\(\psi\) 表示从预测结果到资产权重的映射。

\bookgraphic{0.360}{part-4-67-0001-004-01}{预测模型只是整个金融研究系统中的一个组件。Portfolio Construction 和 Execution Protocol 同样决定最终经济绩效。}

因此，两个 Agent 即使拥有完全相同的预测结果，也可能产生不同的组合收益。

\hypertarget{part-4-67-0001-portfolio-return}{%
\subsubsection{1.2 Portfolio Return}\label{part-4-67-0001-portfolio-return}}

定义资产收益率向量：
\[\boldsymbol r_{t+1}
=
\begin{pmatrix}
r_{1,t+1}\\
\vdots\\
r_{N,t+1}
\end{pmatrix}.\]

在充分投资、忽略交易成本且使用期初权重的条件下，组合的单期简单收益率为：
\[\boxed{
R_{t+1}^{\mathrm{gross}}
=
\boldsymbol w_t^\top\boldsymbol r_{t+1}
}\]

对于只允许做多的 Fully-invested Portfolio：
\[\sum_{i=1}^{N}w_{i,t}=1,
\qquad
w_{i,t}\geq0.\]

如果允许做空、融资或持有现金，则需要额外明确杠杆、现金收益和融资成本的处理方式。

\hypertarget{part-4-67-0001-expected-return-and-portfolio-risk}{%
\subsubsection{1.3 Expected Return and Portfolio Risk}\label{part-4-67-0001-expected-return-and-portfolio-risk}}

定义：
\[\boldsymbol\mu_t
=
\mathbb E[\boldsymbol r_{t+1}\mid\mathcal F_t],\]

\[\Sigma_t
=
\operatorname{Cov}(\boldsymbol r_{t+1}\mid\mathcal F_t).\]

由于合法的投资权重 \(\boldsymbol w_t\) 在时刻 \(t\) 已经确定，即对 \(\mathcal F_t\) 可测，因此：
\[\boxed{
\mathbb E[R_{t+1}^{\mathrm{gross}}\mid\mathcal F_t]
=
\boldsymbol w_t^\top\boldsymbol\mu_t
}\]

以及：
\[\boxed{
\operatorname{Var}(R_{t+1}^{\mathrm{gross}}\mid\mathcal F_t)
=
\boldsymbol w_t^\top\Sigma_t\boldsymbol w_t
}\]

这就是经典 Portfolio Theory 的数学基础。

特别需要注意，组合风险不仅取决于每个资产自身的方差，也取决于资产之间的协方差。

\textbf{Portfolio Risk --- Diversification Explorer}

假想资产 A：预期收益 8\%，波动率 20\%；资产 B：预期收益 5\%，波动率 10\%。各参数采用相同时间尺度。

Weight of Asset A\\
\textbf{50\%}

0.5（范围 0--1）

Correlation \(\rho_{AB}\)\\
\textbf{0.20}

0.2（范围 -1--1）

Expected Portfolio Return

\textbf{6.50\%}

Portfolio Volatility

\textbf{12.04\%}

调整相关系数，可以看到组合的预期收益保持不变，但风险可能发生显著变化。纯数学教学模型。

这说明，即使预测模型的预期收益保持不变，通过改变 Portfolio Construction，也有可能改变风险调整后的表现。

所以，在诊断 RSI 时，我们不能将 Portfolio-level Improvement 全部归因于预测模型变得更准确。

\hypertarget{part-4-67-0001-return-wealth-and-transaction-costs}{%
\subsection{2. Return, Wealth and Transaction Costs}\label{part-4-67-0001-return-wealth-and-transaction-costs}}

\hypertarget{part-4-67-0001-cumulative-returnux7d2fux8ba1ux6536ux76caux7387}{%
\subsubsection{2.1 Cumulative Return（累计收益率）}\label{part-4-67-0001-cumulative-returnux7d2fux8ba1ux6536ux76caux7387}}

沿用 IV-A 中的 Simple Return 定义。

设 \(W_t\) 为时刻 \(t\) 的投资组合财富，\(R_t^{\mathrm{net}}\) 为扣除交易成本后的单期组合收益率。

那么：
\[W_t=W_{t-1}(1+R_t^{\mathrm{net}}).\]

递归展开：
\[\boxed{
W_T=W_0\prod_{t=1}^{T}(1+R_t^{\mathrm{net}})
}\]

相应的累计收益率为：
\[R_{\mathrm{cum}}
=
\prod_{t=1}^{T}(1+R_t^{\mathrm{net}})-1.\]

这里需要区分 Arithmetic Mean Return 与 Geometric Mean Return：
\[\bar R_{\mathrm{arith}}
=
\frac1T\sum_{t=1}^{T}R_t^{\mathrm{net}},\]

\[\bar R_{\mathrm{geo}}
=
\left[
\prod_{t=1}^{T}(1+R_t^{\mathrm{net}})
\right]^{1/T}-1.\]

在财富始终为正的条件下，由 AM--GM Inequality：
\[\boxed{
\bar R_{\mathrm{geo}}
\leq
\bar R_{\mathrm{arith}}
}\]

只有每期收益率完全相同时才达到等号。

这意味着相同的算术平均收益率，可以对应不同的最终财富。

\hypertarget{part-4-67-0001-volatility-dragux6ce2ux52a8ux635fux8017}{%
\subsubsection{2.2 Volatility Drag（波动损耗）}\label{part-4-67-0001-volatility-dragux6ce2ux52a8ux635fux8017}}

考虑两个假想投资组合，它们的算术平均收益率完全相同。

\bookgraphic{0.786}{part-4-67-0001-006-01}{}

Portfolio A --- Final Wealth

\textbf{100.00}

每期收益率 0\%

Portfolio B --- Final Wealth

\textbf{88.47}

交替获得 +20\% 与 −20\%

两者的算术平均收益率均为 0\%，但 Portfolio B 的累计收益为 −11.53\%。合成教学数据。

从数学上看，对于较小的收益率 \(R\)，利用 Taylor Expansion：
\[\log(1+R)\approx R-\frac12R^2.\]

取期望：
\[\mathbb E[\log(1+R)]
\approx
\mu-\frac12(\sigma^2+\mu^2).\]

其中：
\[\mu=\mathbb E[R],
\qquad
\sigma^2=\operatorname{Var}(R).\]

如果 \(\mu^2\) 可以忽略，则：
\[\boxed{
\mathbb E[\log(1+R)]
\approx
\mu-\frac12\sigma^2
}\]

这个近似说明，在其他条件相同的情况下，波动率可能降低长期复利增长率。

不过，这并非对任意收益率分布都成立的精确公式，尤其需要注意大幅收益变化和厚尾分布。

\hypertarget{part-4-67-0001-turnoverux6362ux624bux7387}{%
\subsubsection{2.3 Turnover（换手率）}\label{part-4-67-0001-turnoverux6362ux624bux7387}}

假设上一期投资组合经过市场价格变化之后，其当前权重为 \(\boldsymbol w_t^-\)，而 Agent 希望调整到新的目标权重 \(\boldsymbol w_t\)。

定义双边交易量口径的 Turnover：
\[\boxed{
\mathrm{TO}_t
=
\sum_{i=1}^{N}|w_{i,t}-w_{i,t}^-|
}\]

如果采用单边换手率口径，则通常再除以 2。整个实验必须统一定义。

注意，\(\boldsymbol w_t^-\) 不一定等于上一期目标权重 \(\boldsymbol w_{t-1}\)，因为价格变化本身会改变资产在组合中的相对权重。

假设单位交易金额的比例成本为 \(c\)，则可以构造简化的净收益模型：
\[\boxed{
R_{t+1}^{\mathrm{net}}
\approx
R_{t+1}^{\mathrm{gross}}
-c\cdot\mathrm{TO}_t
}\]

这是忽略成本复合效应、融资费用和市场冲击等因素的一阶近似。

\textbf{Transaction Cost Explorer}

使用 Basis Point（bp）作为收益单位，1 bp = 0.01\%。

Gross return\\
\textbf{12 bp}

12（范围 0--50）

Turnover\\
\textbf{150\%}

1.5（范围 0--3）

Cost per unit traded\\
\textbf{10 bp}

10（范围 0--30）

Gross Return

\textbf{12 bp}

Trading Cost

\textbf{15.0 bp}

Net Return

\textbf{-3.0 bp}

假想成本模型。实际交易成本还可能包括 Bid--Ask Spread、Slippage、Market Impact 和融资成本。

这就解释了为什么某些金融预测模型虽然具有不错的 IC，但其经济价值却非常有限。

对于 RSI，Agent 如果不断增加短周期因子、调整交易频率或者增加组合换手率，就可能出现：
\[R^{\mathrm{gross}}\uparrow,
\qquad
R^{\mathrm{net}}\downarrow.\]

这种现象未必属于 Statistical Overfitting，也可能只是 Objective Function 没有正确计入实际成本。

\hypertarget{part-4-67-0001-sharpe-ratioux590fux666eux6bd4ux7387}{%
\subsection{3. Sharpe Ratio（夏普比率）}\label{part-4-67-0001-sharpe-ratioux590fux666eux6bd4ux7387}}

这是整个 IV-D 最核心的评价指标。

\hypertarget{part-4-67-0001-motivation}{%
\subsubsection{3.1 Motivation}\label{part-4-67-0001-motivation}}

假设我们拥有三个投资组合：

\begin{readertable}
\begin{longtable}[]{@{}
  >{\raggedright\arraybackslash}p{(\columnwidth - 4\tabcolsep) * \real{0.25}}
  >{\raggedleft\arraybackslash}p{(\columnwidth - 4\tabcolsep) * \real{0.44}}
  >{\raggedleft\arraybackslash}p{(\columnwidth - 4\tabcolsep) * \real{0.27}}@{}}
\toprule
Portfolio & Expected Excess Return & Volatility \\
\midrule
\endhead
A & 10\% & 20\% \\
B & 8\% & 8\% \\
C & 15\% & 40\% \\
\bottomrule
\end{longtable}
\end{readertable}

如果仅根据预期收益进行比较，那么 C 表现最好。

但 C 同时承担了更大的风险。

为了衡量单位风险对应的超额收益，我们引入 Sharpe Ratio。

\hypertarget{part-4-67-0001-definition}{%
\subsubsection{3.2 Definition}\label{part-4-67-0001-definition}}

本节新定义 \(Y_t\) 为组合在时刻 \(t\) 的净超额收益率：
\[Y_t=R_t^{\mathrm{net}}-r_{f,t}.\]

设：
\[\mu=\mathbb E[Y_t],\]

\[\sigma^2=\operatorname{Var}(Y_t)>0.\]

Population Sharpe Ratio 定义为：
\[\boxed{
\mathrm{SR}
=
\frac{\mu}{\sigma}
}\]

它衡量单位超额收益波动所对应的平均超额收益。

\textbf{Risk-adjusted Performance Comparison}

\textbf{Portfolio A}

SR = 0.500

Excess return: 10\% · Volatility: 20\%

\textbf{Portfolio B}

SR = 1.000

Excess return: 8\% · Volatility: 8\%

\textbf{Portfolio C}

SR = 0.375

Excess return: 15\% · Volatility: 40\%

三个组合的参数为假想数据，采用相同时间尺度。Sharpe Ratio 不反映所有风险维度。

Portfolio B 虽然没有最高收益，却具有最高的 Risk-adjusted Performance。

但 Sharpe Ratio 将上行和下行波动同等计入风险，因此并不是完整的风险评价体系。

\hypertarget{part-4-67-0001-sharpe-ratio-ux7684-scale-invariance}{%
\subsubsection{3.3 Sharpe Ratio 的 Scale Invariance}\label{part-4-67-0001-sharpe-ratio-ux7684-scale-invariance}}

假设我们将整个超额收益序列乘以一个正常数 \(a>0\)：
\[Y_t'=aY_t.\]

那么：
\[\mathbb E[Y_t']=a\mu,\]

\[\operatorname{Std}(Y_t')=a\sigma.\]

因此：
\[\begin{aligned}
\mathrm{SR}'
&=\frac{a\mu}{a\sigma}\\
&=\frac{\mu}{\sigma}.
\end{aligned}\]

即：
\[\boxed{
\mathrm{SR}(aY)=\mathrm{SR}(Y),
\qquad a>0.
}\]

在理想的线性收益模型下，单纯增加风险敞口不会提高 Sharpe Ratio。

不过，真实市场中的交易成本、融资约束和杠杆导致的非线性损失，可能破坏这种不变性。

\hypertarget{part-4-67-0001-maximum-sharpe-ratio-linear-algebra-perspective}{%
\subsubsection{3.4 Maximum Sharpe Ratio --- Linear Algebra Perspective}\label{part-4-67-0001-maximum-sharpe-ratio-linear-algebra-perspective}}

这一节将 Sharpe Ratio 与你熟悉的 Linear Algebra 联系起来。

设资产的超额收益均值向量为 \(\boldsymbol\mu\)，协方差矩阵为正定矩阵 \(\Sigma\)。

对于非零风险资产敞口 \(\boldsymbol w\)，在允许现金调节、忽略成本和其他仓位约束的理想模型中：
\[\mathrm{SR}(\boldsymbol w)
=
\frac{\boldsymbol w^\top\boldsymbol\mu}
{\sqrt{\boldsymbol w^\top\Sigma\boldsymbol w}}.\]

令：
\[\boldsymbol u=\Sigma^{1/2}\boldsymbol w,
\qquad
\boldsymbol v=\Sigma^{-1/2}\boldsymbol\mu.\]

那么：
\[\mathrm{SR}(\boldsymbol w)
=
\frac{\boldsymbol u^\top\boldsymbol v}
{\|\boldsymbol u\|_2}.\]

由 Cauchy--Schwarz Inequality：
\[\boldsymbol u^\top\boldsymbol v
\leq
\|\boldsymbol u\|_2\|\boldsymbol v\|_2.\]

因此：
\[\mathrm{SR}(\boldsymbol w)
\leq
\|\Sigma^{-1/2}\boldsymbol\mu\|_2.\]

最终得到：
\[\boxed{
\mathrm{SR}_{\max}
=
\sqrt{\boldsymbol\mu^\top\Sigma^{-1}\boldsymbol\mu}
}\]

当 \(\boldsymbol\mu\neq\boldsymbol 0\) 时，最优风险资产敞口方向为：
\[\boxed{
\boldsymbol w^*
\propto
\Sigma^{-1}\boldsymbol\mu
}\]

这里比例系数为正。

\hypertarget{part-4-67-0001-research-implication}{%
\paragraph{Research Implication}\label{part-4-67-0001-research-implication}}

上述推导假设已知真实的 \(\boldsymbol\mu\) 与 \(\Sigma\)。

但现实中只能使用估计量：
\[\hat{\boldsymbol w}
\propto
\hat\Sigma^{-1}\hat{\boldsymbol\mu}.\]

当收益均值估计包含大量噪声，或者 \(\hat\Sigma\) 接近奇异矩阵时，组合权重可能极其不稳定。

因此，即使金融预测模型本身没有发生明显过拟合，\textbf{Portfolio Optimization 也可能独立地引入另一层过拟合。}

这与 IV-C 中 Ridge Regression 对小特征值的处理具有相似的数学结构。

\hypertarget{part-4-67-0001-sharpe-ratio-is-a-statistical-estimator}{%
\subsection{4. Sharpe Ratio Is a Statistical Estimator}\label{part-4-67-0001-sharpe-ratio-is-a-statistical-estimator}}

现在进入本阶段最重要的数学问题。

此前，我们一直在讨论 Population Sharpe Ratio：
\[\mathrm{SR}=\frac{\mu}{\sigma}.\]

但现实中，\(\mu\) 和 \(\sigma\) 都是未知参数。

因此，Agent 实际能够观察到的只是 Empirical Sharpe Ratio：
\[\boxed{
\widehat{\mathrm{SR}}
=
\frac{\hat\mu}{\hat\sigma}
}\]

其中：
\[\hat\mu=\frac1n\sum_{t=1}^{n}Y_t,\]

\[\hat\sigma^2
=
\frac1{n-1}
\sum_{t=1}^{n}(Y_t-\hat\mu)^2.\]

这个统计量同时包含两种估计误差：均值估计误差与方差估计误差。

所以，Sharpe Ratio 本身也具有 Sampling Distribution。

Andrew W. Lo 在 2002 年发表的 \emph{The Statistics of Sharpe Ratios} 中，系统研究了 IID Returns、Stationary Returns 和 Time Aggregation 条件下 Sharpe Ratio 的统计性质。

\href{https://rpc.cfainstitute.org/research/financial-analysts-journal/2002/the-statistics-of-sharpe-ratios}{来源 1}

\hypertarget{part-4-67-0001-mathematical-foundation-delta-method}{%
\subsubsection{4.1 Mathematical Foundation --- Delta Method}\label{part-4-67-0001-mathematical-foundation-delta-method}}

这是一个值得提前掌握的统计工具。

假设估计量 \(\hat{\boldsymbol\theta}_n\) 满足：
\[\sqrt n
(\hat{\boldsymbol\theta}_n-\boldsymbol\theta)
\xrightarrow{d}
\mathcal N(\boldsymbol 0,V).\]

对于在 \(\boldsymbol\theta\) 附近连续可微的函数 \(f\)，通过一阶 Taylor Expansion：
\[f(\hat{\boldsymbol\theta}_n)
\approx
f(\boldsymbol\theta)
+
\nabla f(\boldsymbol\theta)^\top
(\hat{\boldsymbol\theta}_n-\boldsymbol\theta).\]

因此：
\[\boxed{
\sqrt n
\left(
f(\hat{\boldsymbol\theta}_n)
-f(\boldsymbol\theta)
\right)
\xrightarrow{d}
\mathcal N
\left(
0,\nabla f^\top V\nabla f
\right)
}\]

这就是 Multivariate Delta Method。

其思想非常自然：如果我们知道某个估计量的渐近分布，那么可以通过函数的一阶局部线性近似，推导该估计量经过非线性变换后的渐近分布。

\bookgraphic{0.495}{part-4-67-0001-010-01}{Delta Method 的几何直觉：当估计量集中于真实参数附近时，非线性函数可以在局部近似为线性函数。}

这项工具不仅适用于 Sharpe Ratio，也可以用于比率估计量、Correlation Coefficient、MLE 的函数以及其他复杂统计指标。

\hypertarget{part-4-67-0001-asymptotic-distribution-of-sharpe-ratio}{%
\subsection{5. Asymptotic Distribution of Sharpe Ratio}\label{part-4-67-0001-asymptotic-distribution-of-sharpe-ratio}}

\hypertarget{part-4-67-0001-gaussian-case}{%
\subsubsection{5.1 Gaussian Case}\label{part-4-67-0001-gaussian-case}}

假设组合的单期超额收益率满足：
\[Y_t\overset{\mathrm{i.i.d.}}{\sim}
\mathcal N(\mu,\sigma^2).\]

定义：
\[\boldsymbol\theta=
\begin{pmatrix}
\mu\\
\sigma^2
\end{pmatrix},
\qquad
f(\mu,\sigma^2)
=
\frac{\mu}{\sqrt{\sigma^2}}.\]

于是：
\[f(\boldsymbol\theta)=\mathrm{SR}.\]

由 CLT 以及正态总体样本方差的性质：
\[\sqrt n
\begin{pmatrix}
\hat\mu-\mu\\
\hat\sigma^2-\sigma^2
\end{pmatrix}
\xrightarrow{d}
\mathcal N
\left[
\begin{pmatrix}
0\\0
\end{pmatrix},
\begin{pmatrix}
\sigma^2&0\\
0&2\sigma^4
\end{pmatrix}
\right].\]

这里协方差项为零，源于正态总体下样本均值与样本方差的独立性。

\hypertarget{part-4-67-0001-theorem-and-proof}{%
\subsubsection{5.2 Theorem and Proof}\label{part-4-67-0001-theorem-and-proof}}

首先计算梯度：
\[\nabla f
=
\begin{pmatrix}
\frac1\sigma\\[4pt]
-\frac{\mu}{2\sigma^3}
\end{pmatrix}.\]

根据 Delta Method：
\[\begin{aligned}
V_{\mathrm{SR}}
&=
\nabla f^\top V\nabla f\\
&=
\frac1{\sigma^2}\sigma^2
+
\frac{\mu^2}{4\sigma^6}2\sigma^4\\
&=
1+\frac{\mu^2}{2\sigma^2}\\
&=
1+\frac12\mathrm{SR}^2.
\end{aligned}\]

所以：
\[\boxed{
\sqrt n
(\widehat{\mathrm{SR}}-\mathrm{SR})
\xrightarrow{d}
\mathcal N
\left(
0,1+\frac12\mathrm{SR}^2
\right)
}\]

等价地，在大样本条件下：
\[\boxed{
\widehat{\mathrm{SR}}
\overset{\mathrm{approx}}{\sim}
\mathcal N
\left(
\mathrm{SR},
\frac{1+\frac12\mathrm{SR}^2}{n}
\right)
}\]

\(\square\)

因此，其 Standard Error 可以近似估计为：
\[\boxed{
\widehat{\operatorname{SE}}
(\widehat{\mathrm{SR}})
=
\sqrt{
\frac{1+\frac12\widehat{\mathrm{SR}}^2}{n}
}
}\]

这是 Lo (2002) 在 IID Gaussian 条件下得到的经典结果。

\href{https://rpc.cfainstitute.org/research/financial-analysts-journal/2002/the-statistics-of-sharpe-ratios}{来源 1}

\hypertarget{part-4-67-0001-why-this-matters-for-rsi}{%
\subsubsection{5.3 Why This Matters for RSI}\label{part-4-67-0001-why-this-matters-for-rsi}}

假设某个 Agent 的日度 Sharpe Ratio 估计值为：
\[\widehat{\mathrm{SR}}_{\mathrm{daily}}=0.063.\]

如果使用 252 个交易日进行常规年化：
\[\widehat{\mathrm{SR}}_{\mathrm{annual}}
\approx1.\]

但是，估计值为 1，并不意味着真实年化 Sharpe Ratio 就接近 1。

\textbf{Sharpe Ratio --- Estimation Uncertainty}

IID Gaussian 日收益率模型，假设每年有 252 个交易日。

Observed annualized Sharpe Ratio\\
\textbf{1.00}

1（范围 0--3）

Observation length\\
\textbf{1 years}

1（范围 1--20）

\bookgraphic{0.830}{part-4-67-0001-012-01}{}

Standard Error

\textbf{1.001}

Approximate 95\% Confidence Interval

\textbf{{[}-0.96, 2.96{]}}

该置信区间包含零，无法在对应的近似双侧检验下排除真实 Sharpe Ratio 为零。 所有结果均建立在固定策略及相应分布假设下。

sr = 3

\bookgraphic{0.830}{part-4-67-0001-012-state-01}{}

这里有一个相当 Nontrivial 的结果：

\begin{readertable}
\begin{longtable}[]{@{}
  >{\raggedleft\arraybackslash}p{(\columnwidth - 4\tabcolsep) * \real{0.34}}
  >{\raggedleft\arraybackslash}p{(\columnwidth - 4\tabcolsep) * \real{0.34}}
  >{\raggedleft\arraybackslash}p{(\columnwidth - 4\tabcolsep) * \real{0.29}}@{}}
\toprule
Historical Data Length & Observed Annualized SR & Approximate 95\% CI \\
\midrule
\endhead
1 year & 1.00 & {[}−0.96, 2.96{]} \\
5 years & 1.00 & {[}0.12, 1.88{]} \\
10 years & 1.00 & {[}0.38, 1.62{]} \\
20 years & 1.00 & {[}0.56, 1.44{]} \\
\bottomrule
\end{longtable}
\end{readertable}

即使是在相对理想的 IID Gaussian Model 中，Sharpe Ratio 的估计也可能存在相当大的统计不确定性。

而 Financial RSI 可能在同一批数据上反复比较几百个甚至更多候选策略。

这意味着我们不仅面临 Estimation Error，还面临 Selection Bias。

\hypertarget{part-4-67-0001-sharpe-ratio-under-non-normal-returns}{%
\subsection{6. Sharpe Ratio under Non-normal Returns}\label{part-4-67-0001-sharpe-ratio-under-non-normal-returns}}

我们不能忘记 IV-A 中讨论的 Heavy Tails 与 Asymmetric Return Distribution。

Gaussian Assumption 往往不足以描述真实金融收益率。因此，需要将上一节的公式推广。

\hypertarget{part-4-67-0001-skewness-and-kurtosis}{%
\subsubsection{6.1 Skewness and Kurtosis}\label{part-4-67-0001-skewness-and-kurtosis}}

定义收益率的标准化三阶矩：
\[\boxed{
\gamma_3=
\frac{\mathbb E[(Y_t-\mu)^3]}{\sigma^3}
}\]

即 Skewness（偏度）。

以及标准化四阶矩：
\[\boxed{
\gamma_4=
\frac{\mathbb E[(Y_t-\mu)^4]}{\sigma^4}
}\]

即 Kurtosis（峰度）。

注意这里使用的是 Raw Kurtosis，而不是 Excess Kurtosis。

对于 Normal Distribution：
\[\gamma_3=0,\qquad\gamma_4=3.\]

\textbf{Negative Skewness}

左尾较长

\textbf{Symmetric}

左右对称

\textbf{Positive Skewness}

右尾较长

两个收益分布可以具有相同的均值和方差，却具有完全不同的偏度和峰度。因此，它们虽然具有相同的 Population Sharpe Ratio，其估计可靠性却可能不同。

\hypertarget{part-4-67-0001-generalized-asymptotic-variance}{%
\subsubsection{6.2 Generalized Asymptotic Variance}\label{part-4-67-0001-generalized-asymptotic-variance}}

假设 \(Y_t\) 为 IID，且具有有限四阶矩。

此时样本均值与样本方差的渐近协方差矩阵为：
\[V=
\begin{pmatrix}
\sigma^2 & \gamma_3\sigma^3\\
\gamma_3\sigma^3 & (\gamma_4-1)\sigma^4
\end{pmatrix}.\]

继续使用此前定义的梯度：
\[\nabla f=
\begin{pmatrix}
1/\sigma\\
-\mu/(2\sigma^3)
\end{pmatrix}.\]

通过 Delta Method：
\[\begin{aligned}
V_{\mathrm{SR}}
&=
\nabla f^\top V\nabla f\\
&=
1-\frac{\mu}{\sigma}\gamma_3
+\frac{\mu^2}{4\sigma^2}(\gamma_4-1).
\end{aligned}\]

因此：
\[\boxed{
V_{\mathrm{SR}}
=
1-\gamma_3\mathrm{SR}
+
\frac{\gamma_4-1}{4}\mathrm{SR}^2
}\]

最终得到：
\[\boxed{
\operatorname{SE}(\widehat{\mathrm{SR}})
\approx
\sqrt{
\frac{
1-\gamma_3\mathrm{SR}
+\frac{\gamma_4-1}{4}\mathrm{SR}^2
}{n}
}
}\]

实际计算时，可以使用对应的样本估计量代替未知总体参数。

这个结果也是 Bailey 与 López de Prado 后来构造 Probabilistic Sharpe Ratio 的数学基础。

\href{https://papers.ssrn.com/sol3/Delivery.cfm/SSRN_ID3594389_code434076.pdf?abstractid=1821643\&type=2}{来源 1} \href{https://ideas.repec.org/a/rsk/journ4/2223785.html}{来源 2}

\hypertarget{part-4-67-0001-consistency-check}{%
\paragraph{Consistency Check}\label{part-4-67-0001-consistency-check}}

将 Gaussian Distribution 的参数代入：
\[\gamma_3=0,\qquad\gamma_4=3,\]

则：
\[V_{\mathrm{SR}}
=
1+\frac{3-1}{4}\mathrm{SR}^2
=
1+\frac12\mathrm{SR}^2.\]

恰好恢复了 §5 的结果。

\hypertarget{part-4-67-0001-effects-of-skewness-and-kurtosis}{%
\subsubsection{6.3 Effects of Skewness and Kurtosis}\label{part-4-67-0001-effects-of-skewness-and-kurtosis}}

\textbf{Non-normality and Sharpe Ratio Uncertainty}

Population Sharpe Ratio\\
\textbf{0.30}

0.3（范围 0.05--1）

Skewness \(\gamma_3\)\\
\textbf{-1.0}

-1（范围 -2--2）

Kurtosis \(\gamma_4\)\\
\textbf{7.0}

7（范围 2--15）

Gaussian asymptotic variance

\begin{verbatim}
  **1.045**
\end{verbatim}

Non-Gaussian asymptotic variance

\begin{verbatim}
  **1.435**
\end{verbatim}

Standard Error Multiplier

\textbf{1.172×}

Relative Sample Requirement

\textbf{1.373×}

固定真实 SR 时，比较达到相同渐近估计精度所需的样本量。参数受矩不等式 \(\gamma_4\geq\gamma_3^2+1\) 约束。

对于正 Sharpe Ratio：

\begin{itemize}
\tightlist
\item
  Negative Skewness 通常会提高渐近估计方差；
\item
  Higher Kurtosis 会提高渐近估计方差；
\item
  在其他条件相同的情况下，需要更多观测数据才能达到相同统计精度。
\end{itemize}

如果收益率不具有有限四阶矩，那么这一渐近方差公式也不再能够直接使用。

这正是 IV-A 中 Moment Existence 问题在金融绩效评价中的实际体现。

\hypertarget{part-4-67-0001-annualizationux5e74ux5316ux4e0e-temporal-dependence}{%
\subsection{7. Annualization（年化）与 Temporal Dependence}\label{part-4-67-0001-annualizationux5e74ux5316ux4e0e-temporal-dependence}}

\hypertarget{part-4-67-0001-why-do-we-multiply-by-sqrt252}{%
\subsubsection{\texorpdfstring{7.1 Why Do We Multiply by \(\sqrt{252}\)?}{7.1 Why Do We Multiply by \textbackslash sqrt\{252\}?}}\label{part-4-67-0001-why-do-we-multiply-by-sqrt252}}

你可能见过以下金融公式：
\[\mathrm{SR}_{\mathrm{annual}}
=
\sqrt{252}\,
\mathrm{SR}_{\mathrm{daily}}.\]

但为什么是平方根，而不是直接乘以 252？

设 \(q\) 为聚合的时期数量，定义超额收益率的加法聚合量：
\[Y_{1:q}^{\mathrm{sum}}
=
\sum_{t=1}^{q}Y_t.\]

假设每一期具有相同均值、方差，且彼此不相关。

那么：
\[\mathbb E[Y_{1:q}^{\mathrm{sum}}]=q\mu,\]

\[\operatorname{Var}(Y_{1:q}^{\mathrm{sum}})
=q\sigma^2.\]

因此：
\[\begin{aligned}
\mathrm{SR}(q)
&=\frac{q\mu}{\sqrt{q\sigma^2}}\\
&=\sqrt q\frac{\mu}{\sigma}.
\end{aligned}\]

即：
\[\boxed{
\mathrm{SR}(q)=\sqrt q\,\mathrm{SR}(1)
}\]

这里必须强调：该推导基于加法聚合收益，而不是对简单收益率进行精确复利累积。

常见的日度 Sharpe 年化属于这一口径。

\hypertarget{part-4-67-0001-what-if-returns-are-autocorrelated}{%
\subsubsection{7.2 What If Returns Are Autocorrelated?}\label{part-4-67-0001-what-if-returns-are-autocorrelated}}

现在去掉不相关假设。

假设 \(Y_t\) 为弱平稳序列，其自相关函数为：
\[\rho_h=\operatorname{Corr}(Y_t,Y_{t+h}).\]

利用 IV-A 中已经推导过的方差公式：
\[\operatorname{Var}
\left(
\sum_{t=1}^{q}Y_t
\right)
=
\sigma^2
\left[
q+2\sum_{h=1}^{q-1}(q-h)\rho_h
\right].\]

所以：
\[\mathrm{SR}(q)
=
\frac{q\mu}{
\sigma\sqrt{
q+2\sum_{h=1}^{q-1}(q-h)\rho_h
}
}.\]

整理得到：
\[\boxed{
\mathrm{SR}(q)
=
\frac{
\sqrt q\,\mathrm{SR}(1)
}{
\sqrt{
1+2\sum_{h=1}^{q-1}
(1-h/q)\rho_h
}
}
}\]

这就是考虑 Serial Correlation 后的 Sharpe Ratio 时间尺度转换公式。

\href{https://rpc.cfainstitute.org/research/financial-analysts-journal/2002/the-statistics-of-sharpe-ratios}{来源 1}

\textbf{Sharpe Ratio --- Autocorrelation Adjustment}

假设超额收益率遵循平稳 AR(1) 结构，常规年化 SR 固定为 1.50。

Autocorrelation parameter \(\phi\)\\
\textbf{0.30}

0.3（范围 -0.5--0.9）

Naive annualization

\begin{verbatim}
  **1.50**
\end{verbatim}

Autocorrelation-adjusted SR

\begin{verbatim}
  **1.102**
\end{verbatim}

Annualization multiplier

\textbf{11.66}

Naive SR relative bias

\textbf{36.1\%}

根据理论自相关函数计算的聚合调整差异，并非有限样本估计量的实际统计偏差。

对于 AR(1)：
\[\rho_h=\phi^h.\]

当 \(q\) 足够大时：
\[\boxed{
\mathrm{SR}(q)
\approx
\sqrt{
q\frac{1-\phi}{1+\phi}
}\,\mathrm{SR}(1)
}\]

你应该已经看出了它与 IV-A 中 Effective Sample Size 公式的联系。

\hypertarget{part-4-67-0001-actual-empirical-evidence}{%
\subsubsection{7.3 Actual Empirical Evidence}\label{part-4-67-0001-actual-empirical-evidence}}

Lo (2002) 在其共同基金和对冲基金的示例研究中发现，若忽略月度收益率的 Serial Correlation，某个对冲基金的年化 Sharpe Ratio 可能被高估多达 65\%，而正确考虑时间依赖后，基金之间的绩效排名也可能发生显著变化。

\href{https://rpc.cfainstitute.org/research/financial-analysts-journal/2002/the-statistics-of-sharpe-ratios}{来源 1}

这不是说所有金融策略的 Sharpe Ratio 都被高估 65\%，而是说明年化规则本身可能产生具有实际意义的评价偏差。

\hypertarget{part-4-67-0001-rsi-research-implication}{%
\paragraph{RSI Research Implication}\label{part-4-67-0001-rsi-research-implication}}

假设某次 Revision 将预测周期从一天改为二十天，或者引入了具有高度持续性的持仓规则。

这可能同时改变：
\[\mu,\quad\sigma,\quad\rho_h.\]

如果 Evaluator 仍然机械地使用原来的年化公式，就可能无法正确比较两个 Agent 版本。

进一步而言，\textbf{时间尺度转换的修正与 Sharpe Ratio 标准误差的修正是两个不同的问题}。后者仍然需要 HAC、GMM 或适当的 Block Bootstrap 来处理时间依赖，而不能简单认为修正年化公式之后，IID 置信区间就自动有效。

\hypertarget{part-4-67-0001-maximum-drawdownux6700ux5927ux56deux64a4}{%
\subsection{8. Maximum Drawdown（最大回撤）}\label{part-4-67-0001-maximum-drawdownux6700ux5927ux56deux64a4}}

Sharpe Ratio 是一个基于收益率均值和方差的指标。

它并不直接考虑损失发生的顺序。

\hypertarget{part-4-67-0001-definition-1}{%
\subsubsection{8.1 Definition}\label{part-4-67-0001-definition-1}}

沿用 §2 中的组合财富 \(W_t>0\)，定义历史最高财富：
\[M_t=\max_{0\leq s\leq t}W_s.\]

时刻 \(t\) 的 Drawdown 为：
\[\boxed{
D_t=1-\frac{W_t}{M_t}
}\]

Maximum Drawdown 为：
\[\boxed{
\mathrm{MDD}
=
\max_{0\leq t\leq T}D_t
}\]

即：
\[\mathrm{MDD}
=
\max_{0\leq s\leq t\leq T}
\frac{W_s-W_t}{W_s}.\]

\textbf{Maximum Drawdown}

\bookgraphic{0.448}{part-4-67-0001-019-01}{}

\[\mathrm{MDD}=\frac{140-98}{140}=30\%\]

假想财富轨迹。即使最后恢复并创下新高，历史最大回撤仍然为 30\%。

\hypertarget{part-4-67-0001-a-nontrivial-counterexample-same-sharpe-different-risk}{%
\subsubsection{8.2 A Nontrivial Counterexample --- Same Sharpe, Different Risk}\label{part-4-67-0001-a-nontrivial-counterexample-same-sharpe-different-risk}}

假设两个组合拥有完全相同的一组收益率，只是发生顺序不同。

\begin{readertable}
\begin{longtable}[]{@{}
  >{\raggedleft\arraybackslash}p{(\columnwidth - 4\tabcolsep) * \real{0.25}}
  >{\raggedleft\arraybackslash}p{(\columnwidth - 4\tabcolsep) * \real{0.36}}
  >{\raggedleft\arraybackslash}p{(\columnwidth - 4\tabcolsep) * \real{0.36}}@{}}
\toprule
Period & Portfolio A & Portfolio B \\
\midrule
\endhead
1 & +10\% & +10\% \\
2 & −10\% & +10\% \\
3 & +10\% & +10\% \\
4 & −10\% & −10\% \\
5 & +10\% & −10\% \\
6 & −10\% & −10\% \\
\bottomrule
\end{longtable}
\end{readertable}

假设无风险收益率为零。

由于这两个收益率序列的 Multiset 完全相同，因此它们拥有相同的：
\[\hat\mu,\quad\hat\sigma,
\quad\widehat{\mathrm{SR}},
\quad R_{\mathrm{cum}}.\]

但 Maximum Drawdown 并不相同。

\bookgraphic{0.786}{part-4-67-0001-020-01}{}

Portfolio A --- MDD

\textbf{11.79\%}

Portfolio B --- MDD

\textbf{27.10\%}

Both --- Empirical SR

\textbf{0.00}

Both --- Cumulative Return

\textbf{−2.97\%}

这严格说明：
\[\boxed{
\text{Same Sharpe Ratio}
\not\Rightarrow
\text{Same Risk Profile}
}\]

原因在于 Sharpe Ratio 对收益率序列的排列顺序不敏感，而 Maximum Drawdown 是 Path-dependent Statistic。

这也解释了为什么只优化 Sharpe Ratio 可能产生未被该指标充分反映的风险。

\hypertarget{part-4-67-0001-financial-performance-metrics-consolidated-comparison}{%
\subsection{9. Financial Performance Metrics --- Consolidated Comparison}\label{part-4-67-0001-financial-performance-metrics-consolidated-comparison}}

除了 Sharpe Ratio 和 MDD，量化金融中还有其他常见指标。

\begin{readertable}
\begin{longtable}[]{@{}
  >{\raggedright\arraybackslash}p{(\columnwidth - 4\tabcolsep) * \real{0.30}}
  >{\raggedright\arraybackslash}p{(\columnwidth - 4\tabcolsep) * \real{0.36}}
  >{\raggedright\arraybackslash}p{(\columnwidth - 4\tabcolsep) * \real{0.30}}@{}}
\toprule
Metric & Definition & Main Purpose \\
\midrule
\endhead
Cumulative Return & \(\prod_t(1+R_t)-1\) & 总财富增长 \\
Annualized Volatility & 同一口径下的年化标准差 & 收益波动程度 \\
Sharpe Ratio & 超额收益 / 超额收益标准差 & 风险调整绩效 \\
Sortino Ratio & 超目标收益 / 下行偏差 & 强调下行风险 \\
Information Ratio & 主动收益均值 / Tracking Error & 相对于基准的表现 \\
Maximum Drawdown & 最大历史峰谷跌幅 & 路径风险 \\
Calmar Ratio & 年化收益 / MDD & 收益与回撤的关系 \\
Turnover & 投资组合交易量比例 & 交易成本与稳定性 \\
Alpha & 风险模型回归截距 & 风险调整后的异常收益 \\
\bottomrule
\end{longtable}
\end{readertable}

其中 Sharpe Ratio 与 Information Ratio 容易混淆。

如果用 \(R_t^b\) 表示 Benchmark Return，则：
\[\mathrm{IR}
=
\frac{
\mathbb E[R_t-R_t^b]
}{
\operatorname{Std}(R_t-R_t^b)
}.\]

而 Sharpe Ratio 通常以无风险收益率作为比较基准。

因此，二者回答的是不同的金融研究问题。

\hypertarget{part-4-67-0001-why-higher-ic-does-not-guarantee-higher-sharpe-ratio}{%
\subsection{10. Why Higher IC Does Not Guarantee Higher Sharpe Ratio}\label{part-4-67-0001-why-higher-ic-does-not-guarantee-higher-sharpe-ratio}}

现在回到 IV-C 中的 Information Coefficient。

\[\mathrm{IC}_t
=
\operatorname{Corr}_i
(\hat r_{i,t+1},r_{i,t+1}).\]

它衡量预测结果与未来收益之间的横截面相关性。

而：
\[\mathrm{SR}
=
\frac{
\mathbb E[R_t^{\mathrm{net}}-r_{f,t}]
}{
\operatorname{Std}(R_t^{\mathrm{net}}-r_{f,t})
}\]

衡量的是经过组合构建、交易执行和风险暴露之后的经济绩效。

两者之间隔着多个系统组件。

\textbf{Predictive Information} \(\mathrm{IC},\quad\mathrm{RankIC},\quad R^2\) ↓ \textbf{Portfolio Construction}

Ranking · Weighting · Risk exposure

↓ \textbf{Trading Execution}

Turnover · Transaction costs · Market impact

↓ \textbf{Net Economic Performance} \(\mathrm{SR},\quad\mathrm{MDD},\quad R_{\mathrm{cum}}\)

考虑下面的假想单期实验。

\begin{readertable}
\begin{longtable}[]{@{}
  >{\raggedright\arraybackslash}p{(\columnwidth - 4\tabcolsep) * \real{0.50}}
  >{\raggedleft\arraybackslash}p{(\columnwidth - 4\tabcolsep) * \real{0.23}}
  >{\raggedleft\arraybackslash}p{(\columnwidth - 4\tabcolsep) * \real{0.23}}@{}}
\toprule
Metric & Agent A & Agent B \\
\midrule
\endhead
Mean IC & 0.04 & \textbf{0.07} \\
Mean Gross Return & 12 bp & \textbf{16 bp} \\
Turnover & 30\% & 180\% \\
Trading Cost per unit & 10 bp & 10 bp \\
Mean Net Return & \textbf{9 bp} & −2 bp \\
\bottomrule
\end{longtable}
\end{readertable}

Agent B 的预测表现更好，毛收益也更高，但是它产生的组合换手率过高，导致净收益反而为负。

上述均值数据不足以直接计算 Sharpe Ratio，因为尚未给出净收益率的标准差。但它足以说明，更好的 Statistical Prediction 不保证更好的 Economic Performance。

对 X Program \#161 而言，这会带来一个新的研究维度：

\textbf{OOS Decay 发生在预测模型层，还是发生在投资组合及交易执行层？}

如果我们只记录最终 Sharpe Ratio，就很难识别具体原因。

\hypertarget{part-4-67-0001-financial-rsi-ux7684ux591aux5c42ux7ea7ux6027ux80fdux8bcaux65ad}{%
\subsection{11. Financial RSI 的多层级性能诊断}\label{part-4-67-0001-financial-rsi-ux7684ux591aux5c42ux7ea7ux6027ux80fdux8bcaux65ad}}

沿用此前的 Agent Version \(A_k\)，建议为每个版本同时记录三组指标。

\hypertarget{part-4-67-0001-multi-level-performance-trajectory}{%
\subsubsection{11.1 Multi-level Performance Trajectory}\label{part-4-67-0001-multi-level-performance-trajectory}}

\textbf{Level 1 --- Prediction Performance}

\bookgraphic{0.786}{part-4-67-0001-022-01}{}

\textbf{Level 2 --- Trading Activity}

\bookgraphic{0.786}{part-4-67-0001-022-02}{}

\textbf{Level 3 --- Economic Performance}

\bookgraphic{0.786}{part-4-67-0001-022-03}{合成研究示例：预测层 OOS 性能仍然改善，但经济绩效开始退化。换手率增加是值得调查的因素，不是已经确认的因果解释。}

这种情况与普通的 Prediction Overfitting 并不相同。

我们可以利用不同指标的变化模式建立初步诊断。

\begin{readertable}
\begin{longtable}[]{@{}
  >{\raggedright\arraybackslash}p{(\columnwidth - 2\tabcolsep) * \real{0.44}}
  >{\raggedright\arraybackslash}p{(\columnwidth - 2\tabcolsep) * \real{0.52}}@{}}
\toprule
Observed Pattern & Candidate Explanation \\
\midrule
\endhead
IS IC ↑，OOS IC ↓ & 预测模型过拟合 \\
OOS IC ↑，OOS Net SR ↓ & 组合构建或执行问题 \\
Gross SR ↑，Net SR ↓ & 成本或换手率问题 \\
Return ↑，SR ↓ & 风险增长超过收益增长 \\
SR ↑，MDD ↑ & 风险分布或路径性质发生变化 \\
Raw SR ↑，Alpha ↓ & 收益可能来自额外风险暴露 \\
Naive SR ↑，Adjusted SR ↓ & 自相关结构影响评价 \\
\bottomrule
\end{longtable}
\end{readertable}

这些都是候选诊断方向，而不是充分的因果证据。

这也是 III-F 中 Failure Taxonomy 的金融领域细化。

\hypertarget{part-4-67-0001-paired-evaluationux914dux5bf9ux8bc4ux4ef7}{%
\subsubsection{11.2 Paired Evaluation（配对评价）}\label{part-4-67-0001-paired-evaluationux914dux5bf9ux8bc4ux4ef7}}

假设我们希望判断某次 Revision 是否真正改善了经济绩效。

两个相邻版本在相同市场时期产生的收益率为：
\[Y_t^{[k]},
\qquad
Y_t^{[k-1]}.\]

上标 \([k]\) 沿用 IV-C 中表示 Agent 版本的记号。

定义收益差：
\[d_t=
Y_t^{[k]}-Y_t^{[k-1]}.\]

那么：
\[\mathbb E[d_t]
=
\mu_k-\mu_{k-1}.\]

并且：
\[\boxed{
\operatorname{Var}(d_t)
=
\sigma_k^2+\sigma_{k-1}^2
-
2\operatorname{Cov}
(Y_t^{[k]},Y_t^{[k-1]})
}\]

如果两个版本对同一市场冲击具有较强正相关性，那么配对比较可能显著降低收益差的估计误差。

这与 IV-B 中的 Revision-level SNR 直接相关：
\[\mathrm{SNR}_{\mathrm{revision},k}
=
\frac{d_k^2}{\operatorname{Var}(\eta_k)}.\]

通过更合理的实验设计降低 Evaluation Noise，可能比单纯扩大 Agent 的搜索能力更有价值。

不过，检验平均收益差不等于检验 Sharpe Ratio 差异。如果需要比较两个 Sharpe Ratio，应当联合考虑两个版本的均值、方差及协方差，或者对配对收益序列执行适当的 Block Bootstrap。

\hypertarget{part-4-67-0001-research-practice-a-minimal-financial-evaluator}{%
\subsection{12. Research Practice --- A Minimal Financial Evaluator}\label{part-4-67-0001-research-practice-a-minimal-financial-evaluator}}

以下是一个简化的金融绩效评价函数。

它可以作为未来 RSI Framework 中外部 Evaluator 的一个基础组件。

\begin{verbatim}
import numpy as np
from scipy.stats import skew, kurtosis


def evaluate_returns(
    returns,
    risk_free=0.0,
    periods_per_year=252
):
    """
    returns:
        One-dimensional net simple returns.

    risk_free:
        Per-period risk-free return,
        scalar or aligned array.

    Annualization and Sharpe SE assume IID
    observations and additive aggregation.
    """

    r = np.asarray(returns, dtype=float)

    if r.ndim != 1 or len(r) < 4:
        raise ValueError("Insufficient return observations")

    if not np.all(np.isfinite(r)) or np.any(r <= -1):
        raise ValueError("Invalid simple returns")

    y = r - np.asarray(risk_free, dtype=float)

    if y.shape != r.shape or not np.all(np.isfinite(y)):
        raise ValueError("Invalid risk-free returns")

    n = len(y)
    mu = y.mean()
    sigma = y.std(ddof=1)

    if sigma <= 0:
        raise ValueError("Zero excess-return volatility")

    sr = mu / sigma
    sr_annual = np.sqrt(periods_per_year) * sr

    # Include initial wealth in drawdown calculation
    wealth = np.concatenate([
        [1.0],
        np.cumprod(1 + r)
    ])

    peaks = np.maximum.accumulate(wealth)
    drawdown = 1 - wealth / peaks

    # IID Gaussian asymptotic standard error
    se_normal = np.sqrt(
        (1 + 0.5 * sr**2) / n
    )

    # IID non-Gaussian approximation
    gamma3 = skew(y, bias=True)
    gamma4 = kurtosis(
        y,
        fisher=False,
        bias=True
    )

    v_sr = (
        1
        - gamma3 * sr
        + (gamma4 - 1) * sr**2 / 4
    )

    se_non_normal = (
        np.sqrt(v_sr / n)
        if v_sr > 0
        else np.nan
    )

    return {
        "n_observations": n,
        "cumulative_return": wealth[-1] - 1,
        "mean_excess_return": mu,
        "excess_return_volatility": sigma,
        "sharpe_per_period": sr,
        "sharpe_annualized_naive": sr_annual,
        "maximum_drawdown": drawdown.max(),
        "skewness": gamma3,
        "kurtosis": gamma4,
        "sharpe_se_normal_iid": se_normal,
        "sharpe_se_non_normal_iid": se_non_normal
    }
\end{verbatim}

这里刻意将年化 Sharpe 命名为：

\begin{verbatim}
sharpe_annualized_naive
\end{verbatim}

将标准误差命名为：

\begin{verbatim}
sharpe_se_non_normal_iid
\end{verbatim}

原因是我们不能让指标名称掩盖其成立条件。

真正用于科研的 Financial Evaluator 还应当处理时间依赖、交易成本、不同版本的样本对齐、风险暴露，以及多重检验。

该代码目前属于教学参考实现，未经过实际运行测试，也不是课题组的正式评价程序。

\hypertarget{part-4-67-0001-research-hypothesis-sharpe-ratio-overoptimization}{%
\subsection{13. Research Hypothesis --- Sharpe Ratio Overoptimization}\label{part-4-67-0001-research-hypothesis-sharpe-ratio-overoptimization}}

我们现在已经拥有足够的知识提出一个新的候选研究机制。

假设 Agent 在第 \(k\) 轮产生 \(m\) 个候选金融策略：
\[A_k^{(1)},\ldots,A_k^{(m)}.\]

每个候选策略具有真实 Sharpe Ratio：
\[\mathrm{SR}(A_k^{(j)}),\]

但 Agent 只能观察：
\[\widehat{\mathrm{SR}}(A_k^{(j)})
=
\mathrm{SR}(A_k^{(j)})
+
\varepsilon_k^{(j)}.\]

如果采用 Greedy Selection：
\[A_{k+1}
=
\arg\max_{A\in\mathcal C_k\cup\{A_k\}}
\widehat{\mathrm{SR}}(A),\]

那么 Agent 实际上是在优化一个具有抽样误差的统计估计量。

考虑一个极端的 Null Model：所有候选策略的真实 Sharpe Ratio 都为零，并且其估计误差近似为相互独立的 Gaussian Variables。

则经过 \(m\) 次搜索后，最高经验 Sharpe Ratio 的渐近规模约为：
\[\boxed{
\mathbb E\left[
\max_{1\leq j\leq m}
\widehat{\mathrm{SR}}_j
\right]
\approx
\sqrt{\frac{2\log m}{n}}
}\]

这是结合 §5 中零真实 SR 的渐近方差，以及 II-C 中 Extreme Value Approximation 得到的简化结果。

\textbf{Maximum Spurious Sharpe Ratio}

假设全部候选策略的真实 Sharpe Ratio 均为零。

Number of tested strategies\\
\textbf{1,000}

3（范围 1--5）

Historical observation length\\
\textbf{5 years}

5（范围 1--20）

True Population SR

\textbf{0.00}

Approximate Best Observed Annualized SR

\textbf{1.66}

采用 IID Gaussian 和独立候选误差的简化渐近模型。该 leading-order approximation 在候选较少时尤其可能高估实际期望最大值，不代表精确有限样本结果。

这已经非常接近我们下一阶段要研究的 Deflated Sharpe Ratio。

因为一个很高的 Sharpe Ratio 究竟意味着什么，不能脱离以下背景进行解释：
\[\boxed{
\text{Observed Performance}
+
\text{Sample Size}
+
\text{Search History}
}\]

如果不知道这个策略是从多少次尝试中筛选出来的，那么我们对于该策略统计可靠性的判断，就缺少了一项非常重要的信息。

\hypertarget{part-4-67-0001-literature-index}{%
\subsection{14. Literature Index}\label{part-4-67-0001-literature-index}}

\textbf{P0 --- The Statistics of Sharpe Ratios}

Andrew W. Lo · 2002 · Financial Analysts Journal, 58(4), 36--52

本阶段最重要的论文。重点阅读 IID Returns、Non-IID Returns 和 Time Aggregation 部分。

\href{https://rpc.cfainstitute.org/research/financial-analysts-journal/2002/the-statistics-of-sharpe-ratios}{正式出版页面} · \href{https://doi.org/10.2469/faj.v58.n4.2453}{DOI: 10.2469/faj.v58.n4.2453}

\textbf{P0 --- The Sharpe Ratio Efficient Frontier}

David H. Bailey \& Marcos López de Prado · 2012 · Journal of Risk

重点研究非正态收益分布、Sharpe Ratio 估计不确定性及 Probabilistic Sharpe Ratio。

\href{https://ideas.repec.org/a/rsk/journ4/2223785.html}{出版记录} · \href{https://papers.ssrn.com/sol3/papers.cfm?abstract_id=1821643}{SSRN 版本}

两篇文献的出版记录均已核对。建议优先阅读 Lo 的论文，因为它将成为 IV-E 的直接数学基础。

\hypertarget{part-4-67-0001-iv-d-conclusion-index}{%
\subsection{15. IV-D --- Conclusion Index}\label{part-4-67-0001-iv-d-conclusion-index}}

\hypertarget{part-4-67-0001-core-conclusions}{%
\subsubsection{15.1 Core Conclusions}\label{part-4-67-0001-core-conclusions}}

\begin{readertable}
\begin{longtable}[]{@{}
  >{\raggedright\arraybackslash}p{(\columnwidth - 4\tabcolsep) * \real{0.13}}
  >{\raggedright\arraybackslash}p{(\columnwidth - 4\tabcolsep) * \real{0.70}}
  >{\raggedright\arraybackslash}p{(\columnwidth - 4\tabcolsep) * \real{0.14}}@{}}
\toprule
ID & Key Conclusion & Index \\
\midrule
\endhead
D1 & Prediction Model 与 Portfolio Construction 是不同组件 & §1 \\
D2 & 组合风险由整个协方差矩阵决定 & §1.3 \\
D3 & 算术平均收益不等于复利增长率 & §2 \\
D4 & 交易成本可以抵消预测能力提升 & §2.3 \\
D5 & Sharpe Ratio 衡量单位波动对应的超额收益 & §3 \\
D6 & Sharpe Ratio 在理想线性缩放下保持不变 & §3.3 \\
D7 & Portfolio Optimization 可以独立引入过拟合 & §3.4 \\
D8 & Sharpe Ratio 是具有抽样误差的统计估计量 & §4 \\
D9 & Delta Method 可以推导非线性统计量的渐近分布 & §4.1 \\
D10 & Gaussian Sharpe Ratio 具有明确的渐近方差公式 & §5 \\
D11 & Skewness 和 Kurtosis 影响估计可靠性 & §6 \\
D12 & 常规 Sharpe 年化依赖时间聚合及相关性假设 & §7 \\
D13 & Maximum Drawdown 是路径依赖统计量 & §8 \\
D14 & 高 IC 不保证高经济绩效 & §10 \\
D15 & 配对评价可以降低部分版本比较误差 & §11.2 \\
D16 & 大规模策略搜索可以产生虚假高 Sharpe Ratio & §13 \\
\bottomrule
\end{longtable}
\end{readertable}

\hypertarget{part-4-67-0001-core-formula-sheet}{%
\subsubsection{15.2 Core Formula Sheet}\label{part-4-67-0001-core-formula-sheet}}

\textbf{Result 1 --- Portfolio Return and Risk}

\[R_{t+1}
=
\boldsymbol w_t^\top\boldsymbol r_{t+1},\]

\[\operatorname{Var}(R_{t+1}\mid\mathcal F_t)
=
\boldsymbol w_t^\top\Sigma_t\boldsymbol w_t.\]

\textbf{Result 2 --- Sharpe Ratio}

\[\mathrm{SR}=\frac{\mu}{\sigma}.\]

\textbf{Result 3 --- Unconstrained Maximum Sharpe Ratio}

\[\mathrm{SR}_{\max}
=
\sqrt{\boldsymbol\mu^\top\Sigma^{-1}\boldsymbol\mu}.\]

\textbf{Result 4 --- Delta Method}

\[\sqrt n(f(\hat{\boldsymbol\theta})-f(\boldsymbol\theta))
\xrightarrow{d}
\mathcal N(0,\nabla f^\top V\nabla f).\]

\textbf{Result 5 --- Gaussian Sharpe Ratio Standard Error}

\[\boxed{
\operatorname{SE}(\widehat{\mathrm{SR}})
\approx
\sqrt{\frac{1+\mathrm{SR}^2/2}{n}}
}\]

\textbf{Result 6 --- Non-Gaussian Sharpe Ratio Standard Error}

\[\boxed{
\operatorname{SE}(\widehat{\mathrm{SR}})
\approx
\sqrt{
\frac{
1-\gamma_3\mathrm{SR}
+\frac{\gamma_4-1}{4}\mathrm{SR}^2
}{n}
}
}\]

\textbf{Result 7 --- Autocorrelation-adjusted Annualization}

\[\boxed{
\mathrm{SR}(q)
=
\frac{\sqrt q\,\mathrm{SR}(1)}
{\sqrt{
1+2\sum_{h=1}^{q-1}(1-h/q)\rho_h
}}
}\]

\textbf{Result 8 --- Maximum Drawdown}

\[\mathrm{MDD}
=
\max_{0\leq t\leq T}
\left(
1-\frac{W_t}{\max_{s\leq t}W_s}
\right).\]

\hypertarget{part-4-67-0001-proof-index}{%
\subsubsection{15.3 Proof Index}\label{part-4-67-0001-proof-index}}

\begin{readertable}
\begin{longtable}[]{@{}
  >{\raggedright\arraybackslash}p{(\columnwidth - 4\tabcolsep) * \real{0.14}}
  >{\raggedright\arraybackslash}p{(\columnwidth - 4\tabcolsep) * \real{0.63}}
  >{\raggedright\arraybackslash}p{(\columnwidth - 4\tabcolsep) * \real{0.19}}@{}}
\toprule
Proof & Mathematical Result & Location \\
\midrule
\endhead
P1 & Portfolio Conditional Mean and Variance & §1.3 \\
P2 & Arithmetic--Geometric Mean Inequality & §2.1 \\
P3 & Volatility Drag Approximation & §2.2 \\
P4 & Sharpe Ratio Scale Invariance & §3.3 \\
P5 & Maximum Sharpe Ratio via Cauchy--Schwarz & §3.4 \\
P6 & Multivariate Delta Method & §4.1 \\
P7 & Gaussian Sharpe Ratio Asymptotic Variance & §5.2 \\
P8 & Non-Gaussian Sharpe Ratio Asymptotic Variance & §6.2 \\
P9 & Autocorrelation-adjusted Sharpe Ratio & §7.2 \\
P10 & Same Sharpe but Different Drawdown & §8.2 \\
\bottomrule
\end{longtable}
\end{readertable}

\hypertarget{part-4-67-0001-the-most-important-research-insight}{%
\subsection{16. The Most Important Research Insight}\label{part-4-67-0001-the-most-important-research-insight}}

到这里，我们可以进一步修正对于 Financial RSI 的理解。

\textbf{Low Signal-to-Noise Financial Environment}

真实收益信号较弱，估计不确定性较高

↓ \textbf{Noisy Performance Estimator} \(\widehat{\mathrm{SR}}=\hat\mu/\hat\sigma\) ↓ \textbf{Repeated Search \& Selection}

反复生成候选策略、评价并保留最高分版本

↓ \textbf{Selection-induced Optimism}

被保留版本的经验绩效可能系统性偏高

↓ \textbf{Recursive Reinforcement}

Memory、Reflection 与 Self-Modification 可能进一步强化偏差

↓

\[\widehat{\mathrm{SR}}_k^{\mathrm{IS}}\uparrow,\qquad\mathrm{SR}_k^{\mathrm{OOS}}\downarrow\]

\textbf{Potential Generalization Decay}

现在的问题已经不仅仅是：

\begin{quote}
Can the agent discover a profitable strategy?
\end{quote}

而是：

\begin{quote}
Can we distinguish a genuinely improved strategy from the statistical winner of an increasingly powerful search process?
\end{quote}

\textbf{Agent 找到的最优金融策略，可能只是有限历史样本中最幸运的统计估计量。}

当然，这条机制链仍然需要实验验证。目前并未确认 X Program \#161 使用 Sharpe Ratio 作为主要 Objective Function，也未确认其任务一定涉及实际组合交易。因此，本阶段的评价体系应被视为候选研究工具，而不是对课题组现有系统的描述。

\hypertarget{part-4-67-0001-next-iv-e-backtest-overfitting-pbo-deflated-sharpe-ratio}{%
\subsection{Next --- IV-E: Backtest Overfitting, PBO \& Deflated Sharpe Ratio}\label{part-4-67-0001-next-iv-e-backtest-overfitting-pbo-deflated-sharpe-ratio}}

IV-A --- Financial Time Series \& DGP

IV-B --- Low SNR \& Predictability

IV-C --- Financial Factors \& ML Models

IV-D --- Returns, Sharpe Ratio \& Evaluation

IV-E --- Backtest Overfitting, PBO \& DSR

IV-F --- Market Regime \& Robustness

IV-E 将是整个 Part IV 中与 X Program \#161 联系最紧密的一节。

\begin{readertable}
\begin{longtable}[]{@{}
  >{\raggedright\arraybackslash}p{(\columnwidth - 2\tabcolsep) * \real{0.47}}
  >{\raggedright\arraybackslash}p{(\columnwidth - 2\tabcolsep) * \real{0.49}}@{}}
\toprule
Topic & Core Question \\
\midrule
\endhead
Backtest Overfitting & 为什么优秀的历史回测可能没有预测价值？ \\
Probabilistic Sharpe Ratio (PSR) & 观察到的 Sharpe 是否显著超过指定基准？ \\
Deflated Sharpe Ratio (DSR) & 如何考虑大量策略搜索带来的绩效膨胀？ \\
Probability of Backtest Overfitting (PBO) & IS 最优策略在 OOS 中退化的概率如何估计？ \\
CSCV & 如何构造用于 PBO 估计的组合对称交叉验证？ \\
Effective Number of Trials & 如何处理大量相关候选策略？ \\
Adaptive Backtest Overfitting & 传统 PBO 和 DSR 能否直接应用于 RSI？ \\
\bottomrule
\end{longtable}
\end{readertable}

最后一个问题尤其值得深入研究。

传统 Backtest Overfitting 理论主要研究从大量候选策略中进行选择，而 RSI 不仅选择策略，还可能根据历史评价反馈不断改变\textbf{产生候选策略的机制本身}。

这意味着我们即将进入一个非常有研究价值的交叉区域：
\[\boxed{
\text{Backtest Overfitting}
+
\text{Adaptive Data Analysis}
+
\text{Recursive Self-Improvement}
}\]

IV-E 不仅会介绍现有金融统计方法，还会分析这些方法在递归 Agent 环境下可能面临的理论局限。
